\section{Related Work}
\label{sec:related}

\para{Tool use in LLMs.}
\citet{schick2023toolformer} introduced \textsc{Toolformer}, showing that LLMs can learn
\emph{when} and \emph{how} to call APIs in a self-supervised manner.  \citet{qin2023toolllm}
scaled this to 16{,}000 real-world APIs with the \toolbench dataset, demonstrating strong
selection accuracy when models are fine-tuned on diverse function-calling data.  These
works focus on training procedures and end-to-end success rates; they treat selection as a
black box and provide no mechanism for understanding \emph{how} a model matches a query to
a tool.  We complement this line by studying the internal representations that underlie
selection decisions.

\para{Tool selection bias.}
\citet{blankenstein2025biasbust} conduct the most direct behavioral predecessor to our
work.  Across seven LLMs, they show that semantic alignment between a query and a tool's
description is the strongest predictor of tool choice, while repeated exposure to specific
API endpoints during pretraining amplifies selection imbalance.  They also find positional
biases---models tend to prefer earlier-listed tools---and propose a
filter-then-sample mitigation strategy.  Unlike their purely behavioral analysis, we move
inside the model: we test the causal role of description content through controlled
perturbation and use layer-wise probing to identify \emph{where} semantic matching is
computed.

\para{Internal representations for hallucination detection.}
\citet{healy2026internal} demonstrate that the final-layer hidden state of models such as
Qwen-7B and LLaMA-3.1-8B contains enough information to detect hallucinated tool calls
with 72--86\% accuracy using a lightweight MLP classifier.  Their key finding---that mean
pooling over the final layer outperforms token-specific representations---motivates our
use of mean pooling in the probing setup.  However, they consider only the final layer and
focus on hallucination detection rather than on which layer first encodes selection-relevant
information.  We extend this analysis to all 12 layers of \gpttwo and distinguish between
coarse category encoding and fine-grained tool position identification.

\para{Mechanistic interpretability.}
\citet{conmy2023acdc} introduced Automated Circuit Discovery (ACDC) for finding minimal
faithful subgraphs in transformer networks responsible for specific behaviors.  Using edge
attribution patching, \citet{tigges2024llm} show that circuit structures are consistent
across model scale and training stages in the Pythia suite, justifying the use of small
models such as \gpttwo for initial mechanistic analysis.  \citet{merullo2024circuit}
demonstrate 78\% overlap between the Indirect Object Identification (IOI) and Colored
Objects circuits in GPT-2 Medium, suggesting that transformers reuse attention heads for
related relational tasks.  Our probing results suggest tool selection leverages the same
middle-to-late layers involved in other relational computations.  \citet{sharkey2025open}
survey open problems in \MI and identify probe-based analysis as a foundational method
while cautioning that probes detect correlations rather than necessarily causal variables;
we address this caveat by pairing probing with causal perturbation experiments.

\para{Positional biases in LLMs.}
Prior work on the ``lost in the middle'' phenomenon \citep{liu2024lost} shows that LLMs
attend more strongly to the beginning and end of long contexts, which underlies the
primacy bias we observe in tool selection.  Our controlled rotation experiment isolates
this effect and quantifies it independently of semantic content.
